\section{Introduction}
\label{sec:intro}

% NOTE: numbers inside \PH{...} come from claim_audit_report.md (recomputed from raw judge files) and must be bound to numbers.tex.
% NOTE: radiology: 1,073 of 1,200 radiology identifiers point to one image file (claim_audit_report.md, X1). Every radiology-dependent statement carries \rad{}.
% NOTE: the Introduction states 3,500 images through \Nimages; reconsider this figure once the radiology decision is made (2,405 distinct image files at present).

Vision-language models (VLMs) are now proposed as assistants for the interpretation of medical images. A single model can accept an image together with a natural-language instruction and return a free-text description, an answer to a targeted question, or a suggested diagnosis and next step. Two families of models are in use. General-purpose VLMs are trained on broad image--text corpora and are applied to medical images without specific adaptation. Domain-adapted models, for example LLaVA-Med \citep{li2023llavamed}, CheXagent \citep{chen2024chexagent} and HuatuoGPT-Vision \citep{chen2024huatuogptvision}, start from a general architecture and are further trained on biomedical image--text data. The two families differ in training data, in the extent of adaptation and in the clinical areas they target, and it is not established in advance which of them produces the more reliable output for a given imaging domain.

A recurrent failure of both families is hallucination, that is, generated content that is not supported by the input image or that contradicts it \citep{ji2023survey,bai2024hallucination}. In general-domain settings, a hallucination is usually counted as an error of uniform weight. In clinical imaging this convention is inadequate, because the consequence of an error depends on its type. A lesion reported on an image that contains none can lead to unnecessary investigation or treatment. A finding that is present in the image but absent from the report can delay care. An incorrect anatomical location, a misjudged severity, or a plausible but unsupported diagnostic inference each carry a different clinical risk, and a non-committal answer that identifies nothing carries a different risk again. A benchmark that reports a single error rate treats these cases as interchangeable and therefore cannot indicate how harmful the errors of a model are, only how frequent they are.

Prior work has addressed parts of this problem. Med-HallMark and its MediHall Score assign hallucinations to five levels of severity and score model output accordingly, for medical visual question answering and report generation \citep{chen2024medhallmark}. CREST, a risk-weighted score defined over a taxonomy of twelve error types that physicians graded as low, medium or high clinical risk, was developed for chest radiograph report impressions \citep{crest2025medrxiv}. HalluCXR evaluates six models on 856 frontal chest radiographs with an eight-category taxonomy and three severity levels, and reports that normal radiographs drew the most severe hallucinations \citep{wang2026hallucxr}. ClinHallu decomposes the reasoning of a medical model into stages in order to locate where hallucinations arise \citep{yang2026clinhallu}. MedHEval groups hallucinations by cause and evaluates seven mitigation methods on radiological data \citep{chang2025medheval}. CARES assesses several dimensions of trustworthiness, hallucination among them, over a large collection of question--answer pairs spanning many imaging modalities \citep{xia2024cares}. Related studies examine whether models abandon a correct image-based answer under social pressure, and report both that medical fine-tuning can increase this vulnerability \citep{xu2025sycophancy} and that the models which hallucinate least can be the most prone to yield to such pressure \citep{aranya2026grounding}. Severity-aware scoring and taxonomies of hallucination therefore have precedent, and neither is presented here as new in itself.

Three limitations of the existing evidence motivate the present work. First, the benchmarks that grade hallucinations by clinical consequence are concentrated in radiology, and in chest radiography in particular. Dermatology, histopathology and ophthalmology are evaluated in several studies, but these differ in task and in the treatment of error, so that results obtained in one domain cannot be placed beside results obtained in another. Second, a severity-weighted score is usually reported as an absolute number. Such a number has no stable interpretation across modalities, because the base rate of abnormal findings, the typical length of a response and the difficulty of the images all differ, and a value that is low for one modality may be high for another. A comparison of domain-adapted models with general-purpose models in each modality requires a common reference. Third, the evidence on whether domain adaptation reduces or increases hallucination is mixed. The studies cited above point in different directions depending on the task, the models and the pressure applied, and the question is seldom examined with confidence intervals that treat the image, and not the response, as the unit of independent evidence.

This article introduces MedViHall, a cross-modality benchmark for severity-weighted hallucination evaluation of medical VLMs, and uses it to examine these three points. The benchmark contains \Nimages{} images drawn from public datasets in four imaging modalities: chest radiography, dermatology, histopathology and ophthalmology. The images are organised into \Nstrata{} design strata that separate normal images from images with abnormal findings of defined types, so that the evaluation can distinguish a model that invents findings from one that misses them. Each image is presented to each model through up to four query types. The first requests an open-ended report. The second asks a targeted yes/no question. The third asks for a diagnosis and the immediate action it would imply. The fourth, applied to normal images only, states a false premise that asserts an abnormality and asks the model to proceed from it. This protocol elicits the behaviours that matter in practice, namely spontaneous description, answering under a direct question, committing to a diagnosis, and resisting an incorrect clinical premise. Twelve open-weight VLMs, nine general-purpose and three domain-adapted, were evaluated, which yielded \Nresponses{} responses.

Each response is decomposed into atomic claims, and each claim is assigned to one of seven hallucination classes: fabrication, omission, misclassification, severity distortion, anatomical confusion, evasion and clinical reasoning error. Five of these classes correspond to categories already used for chest radiography \citep{wang2026hallucxr}; misclassification and evasion are not among the categories of that scheme. Each class carries a weight that reflects its assumed clinical consequence. The weights were assigned by the authors and are examined by sensitivity analysis rather than presented as clinically elicited values. Three metrics are derived from the labelled claims. The severity-weighted medical hallucination rate (\SMHR) is the mean weight over all claims, so that it increases with the frequency and with the severity of errors. The clinical hallucination rate (\CHR) is the proportion of responses that contain at least one fabrication or omission. The specialisation ratio index (\SRI) divides the \SMHR{} of a model in a given modality by that of the best general-purpose model in the same modality, so that a value above one indicates that the model hallucinates more, in the severity-weighted sense, than the strongest general-purpose reference. All intervals are obtained by resampling images.

The benchmark was used to answer four questions. Does the ranking of the models by \SMHR{} separate the models, once sampling uncertainty over images is taken into account? Do the domain-adapted models hallucinate less than the best general-purpose model in their respective modalities? How consistent are the rankings across modalities? Do inference-time mitigation strategies reduce \SMHR{} when each is compared with an appropriate baseline? The principal findings are as follows.
% NOTE: refresh every figure in the following paragraphs from the audited result files before submission.
The five general-purpose models with the lowest \SMHR{} obtain values between \PH{0.529} and \PH{0.562}. The first three of them are not separated from one another, so that no ordering within this group is supported; the remaining general-purpose models have a higher \SMHR{}. Each of the three domain-adapted models has a higher \SMHR{} than the best general-purpose model in each of the four modalities, with \SRI{} values from \PH{1.31} to \PH{1.85}, and all twelve intervals exclude one \rad{(three of the twelve cells are radiology cells, which depend on the unresolved image issue)}. The \SRI{} of every cell remains above one under equal class weights and under random perturbation of the class weights by up to \PH{30\%}. It is a comparison with the best of nine general-purpose models. It does not isolate the effect of medical adaptation, because the models differ in size, architecture and training data, and one of the three domain-adapted models scores better than at least one general-purpose model in every modality. The values for LLaVA-Med are an upper bound, because about \PH{one third} of its responses begin in mid-sentence as a consequence of a decoding defect.
% NOTE: confirm the LLaVA-Med repair status (rerun or report both numbers) before this sentence is final.
Rankings by \SMHR{} in dermatology, histopathology and ophthalmology are positively correlated across the twelve models (Spearman correlations from \PH{0.66} to \PH{0.81}), but the correlations are weaker among the nine general-purpose models alone, and they therefore reflect in part the separation between the two families; the transfer of rankings to and from radiology could not be assessed \rad{}. The \CHR{} does not order the models as \SMHR{} does, which supports the reporting of both. Contrastive decoding lowered the \SMHR{} of one domain-adapted model, HuatuoGPT-Vision, against both the benchmark baseline and a matched greedy baseline, with and without the radiology images, and it was the only mitigation effect with this property. Retrieval-augmented prompting did not lower \SMHR{} consistently, and it increased the share of responses containing a fabrication or omission for four of the five models examined.

The contributions of this article are the following.
\begin{enumerate}
  \item A seven-class taxonomy of hallucination in medical image interpretation with explicit class weights, adapted from categories in use for chest radiography and extended to the other modalities, and three metrics, \SMHR, \CHR{} and \SRI, that report error frequency, error severity and the position of a model relative to a fixed general-purpose reference in each modality.
  \item A benchmark of \Nimages{} images in four modalities and \Nstrata{} design strata, with a query protocol that includes a false-premise query on normal images, and a two-layer pipeline that extracts claims and assigns them to the taxonomy.
  \item A measurement of the gap between domain-adapted and general-purpose models in each modality, accompanied by image-level bootstrap intervals, tests of which orderings the data support, a sensitivity analysis over the class weights, and an assessment of how far rankings agree across modalities.
  \item A comparison of two inference-time mitigation strategies, retrieval-augmented prompting and contrastive decoding, on five models, in which each strategy is compared with the benchmark baseline and with a matched greedy baseline, so that the effect of the strategy can be separated from the effect of the decoding setting.
  \item An audit of the automated labelling pipeline against annotations produced by six simulated annotators, which supports the detection of the presence of any hallucination at the level of the item and shows that the class-specific labels are not validated. No clinician annotation is included.
  % NOTE: item 5 must be revised if clinician annotation becomes available (Version A); the current wording follows annotation_report.md and annotation_paper_paragraphs.tex (Version B).
\end{enumerate}

The remainder of the article is organised as follows. Section~\ref{sec:related} reviews prior work on hallucination in general and in medical VLMs and on severity-aware evaluation. Section~\ref{sec:benchmark} describes the imaging data, the query protocol, the taxonomy and the metrics. Section~\ref{sec:setup} describes the models, the inference settings and the statistical procedures. Section~\ref{sec:results} reports the results, including the comparison between domain-adapted and general-purpose models, the behaviour of each modality, the stability of rankings and the mitigation experiments. Section~\ref{sec:validation} reports the audit of the labelling pipeline. Section~\ref{sec:discussion} discusses the findings, and Section~\ref{sec:limitations} states the limitations of the study, including the author-assigned class weights, the use of quantised models, the restriction to open-weight models, the dependence of the labels on a single local judge model and the incomplete validation of the class-specific labels.
